\documentclass{szb-solution}

\area{Kalkulus}
\title{Differenciálás II}
\subject{Matematika G1}
\subjectCode{BMETE94BG01}
\date{Utoljára frissítve: \today}
\docno{8}

\begin{document}
\maketitle

\subsection{Implicit függvény deriváltjai}

Legyen
$$
  F(x,y)=x^4y+5y^2x-4=0
  \text.
$$
A parciális deriváltak
$$
  F_x=4x^3y+5y^2,
  \qquad
  F_y=x^4+10xy
  \text.
$$
Ha $F_y\neq0$, akkor az implicit függvény $x$ szerinti deriváltja
$$
  \boxed{
  \frac{\mathrm dy}{\mathrm dx}
  =-\frac{F_x}{F_y}
  =-\frac{4x^3y+5y^2}{x^4+10xy}}
  \text.
$$
Ugyanezt közvetlenül is megkapjuk, mert
$$
  4x^3y+x^4y'+5y^2+10xyy'=0
  \text.
$$

Ha viszont $F_x\neq0$, akkor $x$ tekinthető $y$ függvényének, és
$$
  \boxed{
  \frac{\mathrm dx}{\mathrm dy}
  =-\frac{F_y}{F_x}
  =-\frac{x^4+10xy}{4x^3y+5y^2}}
  \text.
$$
A két derivált ott reciprok egymással, ahol mindkettő értelmezett és nem
nulla.

\subsection{Első és második implicit derivált, érintő}

Az egyenlet
$$
  \ln y+xy=1
$$
csak $y>0$ esetén értelmezett. A $P(1,1)$ pont rajta van a görbén, mert
$\ln1+1\cdot1=1$.

Implicit deriválással
$$
  \frac{y'}y+y+xy'=0
  \text,
$$
tehát
$$
  y'\left(\frac1y+x\right)=-y,
  \qquad
  \boxed{y'=-\frac{y^2}{1+xy}}
  \text.
$$
A második deriválthoz újra deriváljuk az eredeti első derivált egyenletét:
$$
  \frac{y''y-(y')^2}{y^2}+2y'+xy''=0
  \text.
$$
Innen
$$
  y''\left(\frac1y+x\right)
  =\frac{(y')^2}{y^2}-2y'
  \text.
$$
$y'=-y^2/(1+xy)$ behelyettesítésével és egyszerűsítéssel
$$
  \boxed{y''=\frac{y^3(3+2xy)}{(1+xy)^3}}
  \text.
$$
A $P(1,1)$ pontban
$$
  y'(1)=-\frac12,
  \qquad
  y''(1)=\frac58
  \text.
$$
Az érintő egyenlete ezért
$$
  y-1=-\frac12(x-1),
  \qquad
  \boxed{y=-\frac12x+\frac32}
  \text.
$$

\subsection{A kör adott meredekségű érintői}

Az $x^2+y^2=25$ kör implicit deriválásából
$$
  2x+2yy'=0,
  \qquad y'=-\frac{x}{y}
  \text.
$$
A $3/4$ meredekség feltétele
$$
  -\frac{x}{y}=\frac34,
  \qquad y=-\frac43x
  \text.
$$
A kör egyenletébe visszahelyettesítve
$$
  x^2+\frac{16}{9}x^2=25,
  \qquad x=\pm3
  \text,
$$
majd $y=-4x/3$. Tehát
$$
  \boxed{(3,-4)\quad\text{és}\quad(-3,4)}
  \text.
$$

\subsection{Inverz függvény és deriváltja}

Legyen $f(x)=5x^3+x-7$. Mivel
$$
  f'(x)=15x^2+1>0
$$
minden valós $x$-re, $f$ szigorúan növekvő. Továbbá a két végtelenben
$\pm\infty$-hez tart, ezért $\mathbb R$-ről $\mathbb R$-re bijektív, tehát
létezik inverze.

Ha $g=f^{-1}$, akkor $5g(x)^3+g(x)-7=x$. Az inverz deriválási szabálya
szerint
$$
  \boxed{g'(x)=\frac{1}{f'(g(x))}
  =\frac{1}{15g(x)^2+1}}
  \text.
$$
Az inverz Cardano-képlettel explicit alakban is felírható:
\begin{align*}
  g(x)
  &=\sqrt[3]{\frac{x+7}{10}
    +\sqrt{\frac{(x+7)^2}{100}+\frac1{3375}}}\\
  &\quad+\sqrt[3]{\frac{x+7}{10}
    -\sqrt{\frac{(x+7)^2}{100}+\frac1{3375}}}
  \text,
\end{align*}
ahol a valós köbgyököket értjük.

$x_0=1$ esetén $f(1)=5+1-7=-1$, ezért
$$
  \boxed{(f^{-1})'(-1)=\frac1{f'(1)}=\frac1{16}}
  \text.
$$

\subsection{Paraméteres görbe deriváltjai}

Legyen
$$
  x(t)=e^t\cos t,
  \qquad y(t)=e^t\sin t
  \text.
$$
Először a paraméter szerinti deriváltak:
$$
  \dot x=e^t(\cos t-\sin t),
  \qquad
  \dot y=e^t(\sin t+\cos t)
  \text.
$$
Ha $\dot x\neq0$, akkor
$$
  \boxed{
  \frac{\mathrm dy}{\mathrm dx}
  =\frac{\dot y}{\dot x}
  =\frac{\sin t+\cos t}{\cos t-\sin t}}
  \text.
$$
A második deriválthoz
$$
  \ddot x=-2e^t\sin t,
  \qquad
  \ddot y=2e^t\cos t
  \text,
$$
így
\begin{align*}
  \frac{\mathrm d^2y}{\mathrm dx^2}
  &=\frac{\ddot y\dot x-\dot y\ddot x}{\dot x^3}\\
  &=\boxed{\frac{2}{e^t(\cos t-\sin t)^3}}
  \text.
\end{align*}
A $t=\pi/6$ paraméterhez tartozó érintő meredeksége
$$
  \frac{1/2+\sqrt3/2}{\sqrt3/2-1/2}
  =\frac{1+\sqrt3}{\sqrt3-1}
  =\boxed{2+\sqrt3}
  \text.
$$

\subsection{Paraméteresen megadott kör}

A kör paraméterezése
$$
  x(t)=5\cos t,
  \qquad y(t)=5\sin t
  \text.
$$
Ezért
$$
  \frac{\mathrm dy}{\mathrm dx}
  =\frac{5\cos t}{-5\sin t}=-\cot t
  \text.
$$
A $3/4$ meredekséghez $\cot t=-3/4$, vagyis $\tan t=-4/3$. Az ehhez
tartozó egységkörbeli pontok $(3/5,-4/5)$ és $(-3/5,4/5)$, ezért a
$5$ sugarú körön
$$
  \boxed{(3,-4)\quad\text{és}\quad(-3,4)}
  \text.
$$

\subsection{L'Hôpital-szabály alkalmazása}

\begin{enumerate}[a)]
  \item A kifejezés $0/0$ típusú, ezért
        $$
          \boxed{
          \lim_{x\to3}\frac{3^x-x^3}{x-3}
          =\lim_{x\to3}(3^x\ln3-3x^2)
          =27\ln3-27}
          \text.
        $$

  \item $x\to3^+$ esetén mindkét logaritmus $-\infty$-hez tart. Az első
        L'Hôpital-lépés után még $0/0$ alakot kapunk, ezért még egyszer
        alkalmazzuk a szabályt:
        \begin{align*}
          \lim_{x\to3^+}\frac{\ln(x-3)}{\ln(e^x-e^3)}
          &=\lim_{x\to3^+}\frac{e^x-e^3}{e^x(x-3)}\\
          &=\lim_{x\to3^+}\frac{e^x}{e^x(x-3)+e^x}
          =\boxed{1}
          \text.
        \end{align*}

  \item A logaritmus miatt itt $x\to0^+$ értendő. Átírva
        $$
          x\ln x=\frac{\ln x}{1/x}
          \text,
        $$
        ami $(-\infty)/(+\infty)$ típusú. Így
        $$
          \boxed{
          \lim_{x\to0^+}x\ln x
          =\lim_{x\to0^+}\frac{1/x}{-1/x^2}
          =\lim_{x\to0^+}(-x)=0}
          \text.
        $$

  \item Szintén $x\to0^+$ értendő. Legyen $y=x^x$. Ekkor
        $\ln y=x\ln x\to0$, tehát
        $$
          \boxed{\displaystyle\lim_{x\to0^+}x^x=e^0=1}
          \text.
        $$

  \item Közös nevezőre hozva
        $$
          \frac{\cos x}{x}-\frac1{\sin x}
          =\frac{\sin x\cos x-x}{x\sin x}
          \text.
        $$
        Ez $0/0$ típusú. Kétszer alkalmazva a L'Hôpital-szabályt:
        \begin{align*}
          \lim_{x\to0}\frac{\sin x\cos x-x}{x\sin x}
          &=\lim_{x\to0}\frac{\cos2x-1}{\sin x+x\cos x}\\
          &=\lim_{x\to0}\frac{-2\sin2x}{2\cos x-x\sin x}
          =\boxed{0}
          \text.
        \end{align*}
\end{enumerate}

\subsection{Mikor nem segít a L'Hôpital-szabály?}

Az eredeti határérték létezik, mert
$$
  \frac{x^2\sin(1/x)}{\sin x}
  =\frac{x}{\sin x}\,x\sin(1/x)
  \longrightarrow1\cdot0=0
  \text.
$$
A tört $0/0$ típusú, tehát formálisan megpróbálható a L'Hôpital-szabály, de
a deriváltak hányadosa
$$
  \frac{2x\sin(1/x)-\cos(1/x)}{\cos x}
$$
nem rendelkezik határértékkel a nullában a $\cos(1/x)$ oszcillációja miatt.
A szabály elegendő feltétele tehát nem teljesül, és ebből a deriválthányadosból
nem számítható ki a határérték. A helyes eredmény:
$$
  \boxed{\displaystyle
  \lim_{x\to0}\frac{x^2\sin(1/x)}{\sin x}=0}
  \text.
$$

\subsection{Egy valós gyök létezése és egyértelműsége}

Legyen $f(x)=3x^5+15x-2$. Mivel polinom, folytonos, továbbá
$$
  f(0)=-2<0,
  \qquad
  f(1)=16>0
  \text.
$$
A Bolzano-tétel szerint van legalább egy gyöke a $(0,1)$ intervallumban.

Tegyük fel, hogy két különböző valós gyök létezik, $a<b$. Ekkor a
Rolle-tétel szerint volna olyan $\xi\in(a,b)$, amelyre $f'(\xi)=0$. Azonban
$$
  f'(x)=15x^4+15>0
$$
minden valós $x$-re, ami ellentmondás. Így
$$
  \boxed{f\text{-nek pontosan egy valós gyöke van}.}
$$

\clearpage
\subsection{Lokális szélsőértékek}

Legyen
$$
  f(x)=\frac{x^3}{x^2-x-2}
  =\frac{x^3}{(x-2)(x+1)},
  \qquad
  D_f=\mathbb R\setminus\{-1,2\}
  \text.
$$
A derivált
\begin{align*}
  f'(x)
  &=\frac{3x^2(x^2-x-2)-x^3(2x-1)}{(x^2-x-2)^2}\\
  &=\frac{x^2(x^2-2x-6)}{(x-2)^2(x+1)^2}
  \text.
\end{align*}
A kritikus pontok
$$
  x=0,
  \qquad
  x=1\pm\sqrt7
  \text.
$$
A nevező az értelmezési tartományon pozitív, és az $x^2$ tényező nem vált
előjelet. Ezért $f'$ előjele az $x^2-2x-6$ másodfokú tényező előjele:
pozitív a két gyökön kívül, negatív közöttük. Következésképpen
$x=1-\sqrt7$ pontban pozitívról negatívra vált, tehát lokális maximum van;
$x=1+\sqrt7$ pontban negatívról pozitívra vált, tehát lokális minimum van.
$x=0$-ban nincs előjelváltás, így ott nincs szélsőérték.

Mivel a két nemnulla kritikus pontban $x^2=2x+6$, egyszerűsíthetjük a
függvényértékeket, és kapjuk:
$$
  \boxed{
  f(1-\sqrt7)=\frac{20-14\sqrt7}{9}
  \quad\text{(lokális maximum)}}
  \text,
$$
$$
  \boxed{
  f(1+\sqrt7)=\frac{20+14\sqrt7}{9}
  \quad\text{(lokális minimum)}}
  \text.
$$

\end{document}
